\documentclass[11pt]{letter}
\usepackage[margin=1in]{geometry}
\usepackage{amsmath,amssymb}

\signature{Raeez Lorgat\\
Perimeter Institute for Theoretical Physics\\
31 Caroline Street North\\
Waterloo, ON N2L 2Y5, Canada\\
\texttt{rlorgat@perimeterinstitute.ca}}

\date{\today}

\begin{document}

\begin{letter}{The Editors\\
\emph{Duke Mathematical Journal}}

\opening{Dear Editors,}

I am writing to submit the manuscript ``The Seven Faces of the
Collision Residue: A Unified Programme for Chiral Algebras, 3d
Holomorphic-Topological QFT, and Calabi--Yau Quantum Groups'' for
publication in the \emph{Duke Mathematical Journal}.

The paper proves that a single algebraic object, the collision
residue $r(z) = \mathrm{Res}^{\mathrm{coll}}_{0,2}(\Theta_{\mathcal{A}})$
of the universal Maurer--Cartan element of a chirally Koszul chiral
algebra, appears in seven independent mathematical traditions under
seven different names. The seven presentations are:
\begin{enumerate}
\item the genus-zero binary twisting morphism in the bar-cobar
 framework (Ginzburg--Kapranov, Hinich);
\item the $R$-matrix governing meromorphic tensor products of line
 operators (Dimofte--Niu--Py);
\item the classical $r$-matrix of the Poisson vertex algebra
 sigma-model (Khan--Zeng);
\item the generating function of commuting differential operators
 for sphere correlators (Gaiotto--Zeng);
\item the classical limit of the quantum $R$-matrix (Drinfeld);
\item the tensor generating the Sklyanin Poisson bracket
 (Semenov-Tian-Shansky);
\item the generating function of the Gaudin Hamiltonians
 (Feigin--Frenkel--Reshetikhin).
\end{enumerate}

The main theorem establishes a seven-way identification chain
$\mathrm{F1} \Leftrightarrow \mathrm{F2} \Leftrightarrow \cdots
\Leftrightarrow \mathrm{F7}$, with each link proved by an
independent argument and verified by an independent compute engine.
The paper also proves that the agreement is controlled by three
numerical invariants $(p_{\max}, k_{\max}, r_{\max})$ classifying
the standard landscape into four shadow classes $G$, $L$, $C$, $M$,
and that the trichotomy of Hamiltonian operator orders
$\{0, 0, \geq 2\}$ for classes $\{G \cup L,\, C,\, M\}$ explains
the anomalous behaviour of the $\beta\gamma$ system as a stratum
separation phenomenon.

The paper bridges vertex algebra theory, quantum groups, integrable
systems, and holomorphic-topological quantum field theory, making
the \emph{Duke Mathematical Journal} a natural venue given its
tradition of publishing work at the interface of algebra,
geometry, and mathematical physics.

This manuscript has not been published elsewhere and is not under
consideration by any other journal.

\closing{Sincerely,}

\end{letter}
\end{document}
